\subsection{The Fourth Sunday: \latin{Rorate}}\label{sec:fourth-sunday}

In 2026 the Fourth Sunday falls on 20 December, the fourth of the seven days of the great O
antiphons; at its Vespers the Church sings \latin{O clavis David} (gloss: O Key of David). The
station is at the Twelve Apostles. The Mass opens \latin{Rorate}, from Isaiah 45:8, the verse
the Office gave as its versicle at First Vespers of the season:

\begin{quote}
\latin{V. Rorate coeli desuper, et nubes pluant justum. R. Aperiatur terra, et germinet
Salvatorem.}\par
\smallskip
Verse. Drop down, ye heavens, from above, and let the skies pour down the Righteous One.\\
Answer. Let the earth open, and let her bring forth the Saviour.\par
{\smallskip\noindent\small ---Latin: \textit{Breviarium Romanum}, pars hiemalis (Pustet, 1888),
p.~132. English: tr. Bute, \textit{The Roman Breviary} (1908), Winter, p.~213.\par}
\end{quote}

\noindent The verse of the prophet runs in full:

\begin{quote}
Drop down dew, ye heavens, from above, and let the clouds rain the just: let the earth be opened,
and bud forth a saviour: and let justice spring up together: I the Lord have created him.\par
{\smallskip\noindent\small ---Isaias 45:8, Douay-Rheims (Challoner).\par}
\end{quote}

The readings are St Paul's first letter to the Corinthians (4:1--5) and the Gospel of John's
preaching in the desert (Luke 3:1--6), the Gospel already read on the Ember Saturday. The Collect
is the last of the four \latin{Excita} prayers and the most urgent of them:

\begin{quote}
Stir up, O Lord, we pray Thee, Thy power, and come among us, and with great might succour us, that
whereas through our sins and wickedness we are sore let and hindered, Thy bountiful grace and mercy
may speedily help and deliver us; Who livest and reignest with God the Father, in the unity of the
Holy Ghost, one God, world without end.\par
{\smallskip\noindent\small ---Collect of the Fourth Sunday of Advent, tr. Bute, \textit{The Roman
Breviary} (1908), Winter, p.~260.\par}
\end{quote}

Bute's old phrase ``sore let and hindered'' means held back and obstructed.

\subsubsection{Judge not before the time}

The Epistle turns the soul's eye away from its own verdicts. Paul asks the Corinthians to account
him a minister of Christ and a steward of the mysteries of God, from whom it is required that he
be found faithful. Their judgment of him matters little to him, he says, and he does not judge
even himself; though he knows nothing against himself, he is not thereby justified, for the one
who judges him is the Lord. Then comes the verse the Sunday turns on:

\begin{quote}
Therefore, judge not before the time: until the Lord come, who both will bring to light the
hidden things of darkness and will make manifest the counsels of the hearts. And then shall every
man have praise from God.\par
{\smallskip\noindent\small ---1 Corinthians 4:5, Douay-Rheims (Challoner).\par}
\end{quote}

The week before Christmas is a good time to hear this. Advent has asked for examination and
amendment, and the danger at its end is a verdict passed too early: a satisfied account of what
has been done, or a bitter one of what has not. Paul hands both to the Lord who comes. The
faithfulness asked of a steward is to go on doing the work given, and to leave the reckoning to
the one who will make the counsels of the heart manifest.

\subsubsection{Prepare the way}

The Gospel names the emperor, the governor and the tetrarchs, and then the word of God coming to
John in the desert; he goes into all the country about the Jordan preaching the baptism of
penance, as it is written in Isaiah:

\begin{quote}
The voice of one crying in the desert: Prepare ye the way of the Lord, make straight in the
wilderness the paths of our God.

Every valley shall be exalted, and every mountain and hill shall be made low, and the crooked
shall become straight, and the rough ways plain.

And the glory of the Lord shall be revealed, and all flesh together shall see, that the mouth of
the Lord hath spoken.\par
{\smallskip\noindent\small ---Isaias 40:3--5, Douay-Rheims (Challoner). St Luke gives the
prophecy in his own form (Luke 3:4--6).\par}
\end{quote}

\noindent Gregory the Great's twentieth homily (Section~\ref{sec:history}) is read at Matins of
this Sunday. As the \textit{Catena Aurea} gives his words on the cry in the desert:

\begin{quote}
John cries in the desert because he brings the glad tidings of redemption to deserted and
forsaken Judæa, but what he cries is explained in the words, Prepare ye the way of the Lord. For
they who preach true faith and good works, what else do they than prepare the way for the Lord's
coming into the hearts of the hearers, that they might make the paths of God straight, forming pure
thoughts in the mind by the word of good preaching.\par
{\smallskip\noindent\small ---Gregory the Great, \textit{Homiliae in Evangelia}, as the
\textit{Catena Aurea} gives it on Luke 3:4, Oxford tr. (1843), p.~110. The Catena here joins words
of Homilies 7 and 20.\par}
\end{quote}

\noindent Gregory then reads the levelling of the land in the same key:

\begin{quote}
But the crooked places are become straight, when the hearts of the wicked, perverted by a course of
injustice, are directed to the rule of justice. But the rough ways are changed to smooth, when
fierce and savage dispositions by the influence of Divine grace return to gentleness and
meekness.\par
{\smallskip\noindent\small ---Gregory the Great, \textit{Homiliae in Evangelia} 20.6, as the
\textit{Catena Aurea} gives it on Luke 3:5, Oxford tr. (1843), p.~112.\par}
\end{quote}

Both works are within reach of a week's effort, with grace: a wrong set right, a debt paid, an
apology made; a sharp tongue held, a quarrel ended.

The same homily does not let the hearer stay comfortable. On John's words to the crowds, ``O
generation of vipers,'' Gregory says:

\begin{quote}
Let us remark that addressing evil children copying the example of evil parents, the Baptist
calleth them a generation of vipers: in that they were envious at the righteous, and persecuted
them; that they repaid evil for evil; that they hunted out ways of harming their neighbours,---in
all these things following the pattern of carnal parents, the prophet likeneth them to a venomous
brood hatched from a venomous stock.\par
{\smallskip\noindent\small ---Gregory the Great, \textit{Homiliae in Evangelia} 20, tr. Bute,
\textit{The Roman Breviary} (1908), Winter, p.~259.\par}
\end{quote}

\noindent Envy of the good, evil returned for evil, ways sought to harm a neighbor: Gregory's list
is short and plain, and it serves as an examination for the last days of the season.
Section~\ref{sec:confession} takes up the sentence on ``the wrath to come'' that stands just
before it.

\medskip
\noindent\textit{Counsel for the week.} Pray the Collect of the Fourth Sunday each morning; its
``speedily help and deliver us'' is the right prayer for a week that is short and crowded. Read 1
Corinthians 4:1--5 on the Sunday and give the judgment of your Advent to the Lord who comes. Choose
one crooked thing to set straight and one rough thing to make smooth before Christmas, and do
them. If you have not yet gone to confession, go this week (Section~\ref{sec:confession}).
